\section{Local constructions for the numerical comparison}
\label{app:local-numerical}

This appendix supplies the explicit marked semistable construction used
in Section~\ref{sec:log-model}, together with two variants of the
numerical comparison.  First, properly supported intersection squares
allow the specialization argument to be made on a possibly nonproper
family.  We then construct the marked model by monomial modifications.
Finally, we give an alternative algebraic-curve comparison inside a
single irreducible component of the jet locus.

\subsection{Intersection numbers with proper support}

\begin{lemma}[Specialization of a supported square]
\label{lem:supported-square}
Let $q:Y\to C$ be a dominant morphism from an integral normal
$\mathbb Q$-factorial quasi-projective complex threefold to a smooth
connected complex algebraic curve.  Let
$A=\sum_i a_iE_i$ be a rational divisor whose support is proper over
$C$.  For a fiber write its fundamental cycle as
$[Y_c]=\sum_jm_j[S_j]$.  Then the number
\[
 Q_c(A)=\sum_jm_j(A|_{S_j})^2
\]
is independent of $c$.

Here, when $S_j$ is proper, its summand is the usual projective
intersection number.  When $S_j$ is not proper, it is not a component
of $A$, and $A|_{S_j}$ is a rational Cartier divisor with proper curve
support.  Its square means the degree obtained by intersecting that
proper curve cycle with the rational line bundle $A|_{S_j}$.
\end{lemma}

\begin{proof}
The morphism $q$ is flat: its local rings are torsion-free over the
local discrete valuation rings of $C$.  Its fibers have pure dimension
two.  Every vertical $E_i$ is proper, since the support of $A$ is
proper over $C$.  A horizontal $E_i$ is projective and flat over $C$:
properness together with quasi-projectivity gives projectivity, and
flatness again follows from torsion-freeness.  All intersections below
therefore have proper support.  Multiplying the divisors by common
positive integers makes them Cartier; we divide by these integers
after taking intersection numbers.

Commutativity and additivity of Cartier-divisor intersections give
\begin{align*}
 Q_c(A)
 &=\sum_i a_i\sum_jm_j
       \deg\bigl(A\cdot E_i\cdot[S_j]\bigr)\\
 &=\sum_i a_i\deg\bigl(A|_{E_i}\cdot Y_c|_{E_i}\bigr).
\end{align*}
For distinct $E_i$ and $S_j$, the two orders of restriction describe
the same proper intersection curves with their intersection
multiplicities.  If $E_i=S_j$, all calculations take place on a proper
surface and use its normal-bundle class.  These observations also
justify the formula when some fiber component is not proper.

For a vertical $E_i$, the line bundle of $Y_c$ restricts trivially to
$E_i$, since it comes from a point divisor on $C$ and $q|_{E_i}$ is
constant.  Its term is zero, including when $E_i\subset Y_c$.
For a horizontal $E_i$, the term is the degree of the fixed rational
line bundle $A|_{E_i}$ on the scheme-theoretic fiber $(E_i)_c$.
That degree is constant in a projective flat curve family.  For
example, it is the coefficient of $n$ in the fiberwise Euler
characteristic of a Cartier multiple of $nA|_{E_i}$; flatness makes
the Euler characteristic independent of $c$, and the intersection
formula identifies this coefficient.  This proves the assertion.
\end{proof}

There is a second proof that explains why a compactification of the
ambient family causes no extra term.  Take an integral normal
projective compactification $\overline Y\to C$ of $Y\to C$.
The support of $A$ is already closed in $\overline Y$: its map to
$\overline Y$ is proper over the separated base $C$.  A Cartier
multiple of $A$ thus extends to a Cartier divisor on $\overline Y$
by gluing it to the zero divisor on
$\overline Y\setminus\operatorname{Supp}(A)$.  The compactification
is flat over $C$.  Lemma~\ref{lem:square-specialization} makes the
extended fiber square constant.  Its restriction to every component
contained in $\overline Y\setminus Y$ is trivial, and therefore has
square zero.  On a component meeting $Y$, its square is the supported
intersection number in Lemma~\ref{lem:supported-square}.  The
projective identity is consequently precisely the desired local one.

\begin{corollary}[The local exceptional-square comparison]
\label{cor:local-exceptional-square}
Suppose a projective birational morphism $h:Y\to X$ is a log model
over a smooth curve as in Section~\ref{sec:full-boundary}, except
that $X\to C$ need not be proper.  Assume $X$ and $Y$ are
quasi-projective, all exceptional prime divisors map into the section
$\Gamma\simeq C$, and put
\[
 A=K_Y+H-h^*K_X.
\]
If $Y_0=S_0+\sum_aS_a$ and the nearby fibers are irreducible, then
\[
 (A|_{Y_c})^2=(A|_{S_0})^2+\sum_a(A|_{S_a})^2.
\]
Under the full-boundary, adjunction, and nefness hypotheses of
Proposition~\ref{prop:beta-from-log-model}, the first two squares
are respectively $-\beta(X_c,\Gamma(c))$ and
$-\beta(X_0,\Gamma(0))$, while
\[
 (A|_{S_a})^2=(K_{S_a}+B_{S_a})^2.
\]
Thus vanishing of the extra squares proves constancy of the local
invariant without a projective compactification of $X$.
\end{corollary}

\begin{proof}
The divisor $A$ is exceptional and has proper support over the section,
because $h$ is projective.  Lemma~\ref{lem:supported-square} applies.
On each $S_a$, a divisor pulled back from $X$ has trivial line-bundle
class, giving the displayed adjunction formula.  On $S_0$ and on a
nearby fiber, pull back to a good resolution of the corresponding
surface germ.  As in Lemma~\ref{lem:projective-local-square}, log
canonicity makes the difference from the pulled-back logarithmic
canonical divisor an effective exceptional negative part, orthogonal
to the relatively nef part.  The pullback of $A$ is the numerical
exceptional representative of that nef part.  Its properly supported
square is therefore $-\beta$, by the definition of the local invariant.
This argument uses no global canonical square of $X_c$.
\end{proof}

\subsection{An explicit marked semistable construction}

We spell out the dimension-three toroidal construction underlying
the marked semistable reduction used in Proposition~\ref{prop:log-model};
compare \cite[Theorems~7.17 and~7.19]{KM1998}.  It preserves the
horizontal markings throughout the construction.

\begin{proposition}[Marked monomial reduction in dimension three]
\label{prop:marked-monomial-reduction}
Let $q:Z\to C$ be a morphism from a smooth quasi-projective complex
threefold to a smooth curve.  Let $0\in C$, write
$Z_0=\sum_i a_iZ_i$, and let $H$ be a reduced divisor with horizontal
components.  Suppose that $(Z_0)_{\mathrm{red}}+H$ is simple normal
crossing, and that away from $0$ the morphism is smooth with relative
simple normal crossing divisor $H$.
After shrinking $C$ about $0$, choose a parameter $\tau$ with no
other zero and an integer $b$ divisible by every $a_i$.  After the
finite base change $s^b=\tau$ and normalization, there is a projective
birational modification whose total space is smooth and whose special
fiber is reduced and simple normal crossing together with the
transformed horizontal marking.  The modification is an isomorphism
on every open set where the original morphism was smooth and the
marking relatively simple normal crossing.
\end{proposition}

\begin{proof}
Work componentwise if the normalized base change is disconnected.
At a point of $Z_0$, choose simple normal crossing coordinates in
which
\[
 \tau=u\prod_{i\in I}z_i^{a_i},
 \qquad u\text{ a unit},\qquad |I|\le3.
\]
The horizontal components present are among the other coordinate
hyperplanes.  An \'etale root of $u$ reduces the local calculation
to the monomial equation.  Each normalized local sheet is toric in
the vertical coordinates, with cone and lattice
\[
 \sigma_I=\mathbb R_{\ge0}^{I},\qquad
 N_I=\left\{(d_i)\in\mathbb Z^{I}:
                \sum_{i\in I}a_id_i\equiv0\pmod b\right\}.
\]
The order of the new base parameter is the integral height function
$\ell(d)=\sum_i a_id_i/b$.  Since $a_i\mid b$, the primitive
generator $(b/a_i)e_i$ of each original ray has height one.
The remaining coordinates, including the horizontal markings, are
unchanged product factors.

Subdivide the cones by inserting all their height-one lattice rays.
There are finitely many: their intersections with $\ell=1$ lie in
compact simplices.  At each stage insert a point not already a ray
by a star subdivision of its smallest containing cone.  Starting
with simplicial cones, this keeps every cone simplicial.  To make
the choices compatible across all charts, carry out the sequence on
the single orthant with one coordinate for each global $Z_i$ and
the congruence lattice with the same coefficients $a_i$.  Restrict
its subdivisions to each face that actually occurs.  A point not
on such a face does not change that face.

Here is a global realization by projective modifications.  A vertical
ray $v$ records the orders along its divisor of the pullbacks of the
original $Z_i$.  Let $G_v$ be the sum of all vertical prime divisors
with that order vector.  On every current chart the fan is simplicial,
so $G_v$ is rational Cartier.  Suppose a ray to be inserted is in
the relative interior of the cone generated by $v_1,\ldots,v_k$,
and write $v=\sum_j\lambda_jv_j$, with all $\lambda_j>0$ rational.
Choose positive integers $c_j$ such that the $c_jG_{v_j}$ are
Cartier and all products $c_j\lambda_j$ are equal.  Normalize the
blowup of the coherent ideal
\[
 \mathcal I=\sum_j\mathcal O(-c_jG_{v_j}).
\]
On a toric chart its generators are monomials.  The regions on which
one of their orders is smallest give exactly the star subdivision
at $v$: the equality of these orders is
$c_1d_1=\cdots=c_kd_k$ in ray coordinates.  On a chart missing one
of the required rays the corresponding summand is the unit ideal,
so the modification is the identity.  Thus these normalized blowups
realize the compatible subdivisions globally.  They are projective,
and preserve the product coordinates for the horizontal markings.

We verify regularity, including the small-dimensional lattice step.
In an actual chart the vertical cone has dimension at most three.
Its section at height one is now a segment, a triangle, or a point
with no lattice points other than its vertices.  An empty lattice
segment is primitive.  An empty lattice triangle is unimodular:
translate one vertex to zero and denote its edge vectors by $v,w$.
If their lattice index exceeded one, a nonzero lattice point of the
half-open parallelogram would have the form
$\alpha v+\gamma w$, with $0\le\alpha,\gamma<1$.  If
$\alpha+\gamma\le1$, this point lies in the triangle and is not a
vertex.  If $\alpha+\gamma>1$, the lattice point
$(1-\alpha)v+(1-\gamma)w$ does so.  Both contradict emptiness.
The edge vectors therefore form a basis of the height-zero lattice.
Adjoining one height-one vertex gives a basis of the full lattice,
so the corresponding cone is regular.  The same conclusion follows
from the primitive-segment calculation in dimension two and is
immediate in dimension one.

The resulting local toric varieties are smooth.  Every vertical ray
still has height one, so the divisor of $s$ is reduced.  Regular
toric coordinates show that this divisor together with all horizontal
coordinate divisors is simple normal crossing.  The original global
labels and order vectors distinguish local boundary branches, so
the construction does not identify distinct branches of one component.
Outside the special fiber the cover is \'etale and no subdivision
acts.  At a smooth point of the original morphism over $0$ there is
only one vertical coordinate of multiplicity one; its cone needs no
subdivision.  This proves the asserted isomorphism property as well.
\end{proof}

\subsection{Algebraic curves through prescribed jets}

The constructibility lemma in Section~\ref{sec:algebraic-reduction}
is useful in its own right.  The scalar comparison also admits the
following argument that does not use constructibility of the scalar.

\begin{lemma}[Comparison within an irreducible jet component]
\label{lem:jet-component-comparison}
Let $U_\mu$ be the locally closed algebraic jet locus of
Lemma~\ref{lem:mu-jet-locus}.  Assume the algebraic constancy assertion
in Proposition~\ref{prop:scalar-curve-reduction}.  Any two points of
the same irreducible component of $U_\mu$, one of which defines a
non-log-canonical surface germ, have the same Wahl invariant.
Consequently that proposition holds without invoking
constructibility of the Wahl invariant.
\end{lemma}

\begin{proof}
Write $Z$ for the irreducible component and let $a_0,a_1\in Z$ be
the points in question.  If they coincide there is nothing to prove.
Resolve a projective closure of $Z$, obtaining a smooth irreducible
projective variety $\widetilde Z$ with a proper surjective map to
that closure.  Choose a point above each $a_i$.  If $\dim Z=1$,
use $\widetilde Z$ itself.  If $\dim Z\ge2$, general sufficiently
ample hypersurfaces through the two chosen points cut out a smooth
irreducible projective curve through them.  Indeed the systems can
be chosen to separate first jets at those points; their general
members are smooth there, and Bertini applies away from the finite
base locus.  Iterating the irreducibility and smoothness assertions
of Bertini gives the required complete-intersection curve; compare
\cite{KleimanBertini}.

The complement of the inverse image of $Z$ is closed and misses the
chosen points.  Its intersection with the curve is therefore
finite.  Delete those finitely many points, and, if needed, restrict
the curve further to an affine open containing the two chosen
points.  The result is a smooth connected algebraic curve $C'$,
with a morphism to $U_\mu$ whose image contains $a_0$ and $a_1$.
Pulling back the universal polynomial gives regular coefficients
on $C'$ and constant Milnor number $\mu$.  Log canonicity is
constant by Lemma~\ref{lem:lc-persistence}; hence all its germs are
non-log-canonical.  Apply the assumed scalar assertion with any
point of $C'$ as the marked point.  Its value agrees with the
nearby general values.  The exceptional algebraic subset of a
curve is finite; shrinking an analytic disk about the marked
point avoids its other points.  The value is therefore constant
on that disk, including its center.  It follows that the invariant
is locally constant throughout the connected curve $C'$, and in
particular has the same value at the two prescribed points.

Finally, the finitely many irreducible components of $U_\mu$ not
containing $a_0$ form a closed subset missing $a_0$.  A sufficiently
small analytic neighborhood of $a_0$ avoids that subset.  Every
nearby jet $a(t)$ in the holomorphic jet map therefore shares an
irreducible component with $a_0$.  The comparison just proved makes
its Wahl invariant equal to the central value.  Individual finite
determinacy then returns this equality to the original holomorphic
family, exactly as in Proposition~\ref{prop:scalar-curve-reduction}.
\end{proof}
